\chapter{Tag 6 — Polynome über GF($2^3$) und die Reed-Solomon-Idee}

\begin{quote}
Gestern hast du GF($2^3$) gebaut — einen Körper mit acht Elementen,
in dem du addieren, multiplizieren und auch dividieren kannst. Heute
nutzen wir diesen Körper, um \textbf{Polynome darüber} zu betrachten.
Genauer: Polynome, deren Koeffizienten GF($2^3$)-Elemente sind. Daraus
ergibt sich — fast wie von selbst — die zentrale Idee von
\textbf{Reed-Solomon}: ein Verfahren, das aus weniger als der Hälfte
aller übertragenen Werte das Original rekonstruieren kann.
\end{quote}

\section*{Lernziele für heute}

Am Ende von Tag 6 kannst du \dots

\begin{itemize}
\item Polynome mit Koeffizienten aus GF($2^3$) addieren, multiplizieren
  und an einer Stelle auswerten
\item das Theorem zur Stützstellen-Eindeutigkeit nennen und an einem
  konkreten Beispiel anwenden („zwei Punkte $\to$ Gerade“)
\item die \textbf{Reed-Solomon-Idee} erklären: Daten als Polynom
  interpretieren, an mehr Stellen auswerten, als für die Rekonstruktion
  nötig sind
\item ein RS$(n, k)$-Codewort von Hand berechnen (für ein kleines
  Beispiel über GF($2^3$))
\item in Python eine \texttt{Polynom}-Klasse über GF($2^3$) bauen und
  damit einen einfachen RS-Encoder schreiben
\end{itemize}

Material: Bleistift, kariertes Papier, deine GF($2^3$)-Multiplikationstafel
aus Tag 5 (Bleistiftübung 6 — falls du sie noch nicht hast, ab und zu
nachschlagen!), Thonny.

% =============================================================
\section{Polynome — diesmal mit Koeffizienten aus GF($2^3$) \normalfont(\textit{$\approx$ 25 min})}
% =============================================================

\subsection*{Erst eine Notations-Sache}

Bei Tag 5 haben wir die GF($2^3$)-Elemente als Polynome geschrieben:
$x^2 + x + 1$, $x + 1$ und so weiter. Das war eine \emph{Implementierungs-Sicht}
— so haben wir den Körper überhaupt erst gebaut.

Ab heute wechseln wir die Brille: \textbf{ein GF($2^3$)-Element ist
einfach eine Zahl zwischen 0 und 7 (oder eine 3-Bit-Folge), mit den
besonderen Rechenregeln aus Tag 5.} Die Polynom-Schreibweise von
gestern brauchen wir nur noch ab und zu zur Erinnerung; wir schreiben
\texttt{3} statt $x+1$, \texttt{5} statt $x^2+1$, \texttt{0} bleibt $0$.

Warum die Umstellung? Weil wir heute \emph{neue} Polynome bauen, deren
Variable \emph{unabhängig} von der GF-internen Notation ist. Damit wir
nicht zweimal mit „$x$“ durcheinanderkommen.

\subsection*{Polynome über GF($2^3$): die Definition}

Ein Polynom über GF($2^3$) ist ein Ausdruck der Form
\[
  f(x) \;=\; a_0 + a_1 x + a_2 x^2 + \dots + a_{k-1} x^{k-1}
\]
wobei jeder \textbf{Koeffizient $a_i$ ein Element von GF($2^3$)} ist
— also eine Zahl aus $\{0, 1, 2, 3, 4, 5, 6, 7\}$.

\textbf{Achtung:} das $x$ ist hier nichts Konkretes — es ist nur ein
Platzhalter. Beim Auswerten setzt du für $x$ irgendein
GF($2^3$)-Element ein, und alle Rechnungen ($+$, $\cdot$) finden in
GF($2^3$) statt.

\subsection*{Wie addiert man zwei solche Polynome?}

Genauso wie in der Schule: koeffizientenweise. Nur dass die Addition
in GF($2^3$) eben \textbf{XOR} ist.

\textit{Beispiel:}
\begin{align*}
  f(x) &= 3 + 5x + 2x^2 \\
  g(x) &= 1 + 7x + 2x^2 \\
  f(x) + g(x) &= (3 \oplus 1) + (5 \oplus 7) \cdot x + (2 \oplus 2) \cdot x^2 \\
              &= 2 + 2x + 0 = 2 + 2x
\end{align*}

\subsection*{Wie multipliziert man?}

Wieder wie in der Schule: jeder Term mit jedem, dann die
gleichgradigen Terme zusammenfassen. Nur eben \emph{mit
GF($2^3$)-Arithmetik} statt der gewohnten Schul-Arithmetik.

\textit{Beispiel:}
\begin{align*}
  (1 + 2x) \cdot (1 + 3x)
    &= 1 \cdot 1 + 1 \cdot 3x + 2x \cdot 1 + 2x \cdot 3x \\
    &= 1 + 3x + 2x + (2 \cdot 3) x^2 \\
    &= 1 + (3 \oplus 2) x + 6 x^2 \\
    &= 1 + x + 6 x^2
\end{align*}

(Aus der Multiplikationstafel von Tag 5: $2 \cdot 3 = 6$, denn
$x \cdot (x+1) = x^2 + x = 6$.)

\begin{bleistiftuebung}{1}{Polynome rechnen}
Halte deine GF($2^3$)-Multiplikationstafel aus Tag 5 bereit.

\begin{enumerate}[label=(\alph*), leftmargin=*]
\item $f(x) = 4 + 3x + x^2$, $g(x) = 5 + 6x$. Berechne $f(x) + g(x)$.
\item Berechne $(1 + x) \cdot (1 + x)$ — also $(1 + x)^2$ — über
  GF($2^3$). Vergleiche mit der „normalen“ Algebra.
\item Berechne $(2 + x) \cdot (3 + x^2)$ und bringe das Ergebnis in die
  Standardform $a_0 + a_1 x + a_2 x^2 + a_3 x^3$.
\end{enumerate}
\end{bleistiftuebung}

% =============================================================
\section{Polynome auswerten — und die Magie der Stützstellen \normalfont(\textit{$\approx$ 30 min})}
% =============================================================

\subsection*{Auswerten}

\textbf{Auswerten} heißt: setze für $x$ ein konkretes
GF($2^3$)-Element ein und rechne den ganzen Ausdruck aus.

\textit{Beispiel:} $f(x) = 3 + 5x + 2x^2$, ausgewertet an $x = 2$:
\begin{align*}
  f(2) &= 3 + 5 \cdot 2 + 2 \cdot 2^2 \\
       &= 3 + 1 + 2 \cdot 4 \quad\quad (\text{denn } 5 \cdot 2 = 1, \; 2 \cdot 2 = 4) \\
       &= 3 + 1 + 3 \quad\quad\quad\;\, (\text{denn } 2 \cdot 4 = 3) \\
       &= 3 \oplus 1 \oplus 3 = 1
\end{align*}

\subsection*{Das zentrale Theorem}

Hier ist die mathematische Aussage, die alles weitere trägt:

\begin{quote}
\textbf{Ein Polynom vom Grad kleiner als $k$ ist eindeutig bestimmt
durch $k$ verschiedene Auswertungen.}
\end{quote}

Anschaulich:
\begin{itemize}
\item Eine Gerade ($f(x) = a_0 + a_1 x$, also Grad $< 2$) ist durch
  \textbf{2} Punkte eindeutig festgelegt.
\item Eine Parabel (Grad $< 3$) ist durch \textbf{3} Punkte eindeutig
  festgelegt.
\item Allgemein: Polynom Grad $< k$ $\to$ $k$ Punkte reichen zur
  vollständigen Rekonstruktion.
\end{itemize}

Das gilt nicht nur für reelle Zahlen, das gilt für \textbf{jeden
Körper} — also auch für GF($2^3$).

\subsection*{Warum ist das wahr?}

(Skizze, kein vollständiger Beweis.) Angenommen, es gäbe zwei
verschiedene Polynome $f$ und $g$, beide vom Grad $< k$, die an $k$
Stellen $x_1, \dots, x_k$ \emph{dieselben} Werte annehmen. Dann hätte
das Differenzpolynom $h(x) = f(x) - g(x)$ an diesen $k$ Stellen
jeweils den Wert $0$ — also $k$ Nullstellen. Aber: ein Polynom vom
Grad $< k$, das nicht das Nullpolynom ist, kann höchstens $k - 1$
Nullstellen haben. Widerspruch! Also muss $h \equiv 0$, also $f = g$.

\begin{bleistiftuebung}{2}{Gerade durch zwei Punkte}
Du hast zwei Auswertungen einer unbekannten Geraden $f(x) = a_0 + a_1 x$
über GF($2^3$):
\[
  f(1) = 4, \qquad f(2) = 7.
\]

\begin{enumerate}[label=(\alph*), leftmargin=*]
\item Stelle das Gleichungssystem für $a_0$ und $a_1$ auf. (Zwei
  Gleichungen, zwei Unbekannte.)
\item Löse es \emph{in GF($2^3$)}. Tipp: Subtraktion ist Addition
  ($a - b = a + b$), und Division ist Multiplikation mit dem
  Inversen.
\item Welches Polynom ist es? Verifiziere durch Auswertung an $x = 1$
  und $x = 2$.
\end{enumerate}
\end{bleistiftuebung}

% =============================================================
\section{Die Reed-Solomon-Idee \normalfont(\textit{$\approx$ 30 min})}
% =============================================================

Jetzt wird's spannend. Die ganze Vorbereitung der letzten zwei Tage
zahlt sich aus.

\subsection*{Das Spiel}

\begin{enumerate}
\item Du hast $k$ \textbf{Datensymbole} $a_0, a_1, \dots, a_{k-1}$,
  jedes ein GF($2^3$)-Element. Interpretiere sie als die Koeffizienten
  eines Polynoms:
  \[
    f(x) = a_0 + a_1 x + a_2 x^2 + \dots + a_{k-1} x^{k-1}.
  \]
\item Wähle $n > k$ \textbf{verschiedene Stützstellen}
  $x_1, x_2, \dots, x_n$ in GF($2^3$). (Bei GF($2^3$) gibt es nur 8
  Elemente, also höchstens $n = 8$.)
\item \textbf{Sende:} die $n$ Auswertungen $f(x_1), f(x_2), \dots,
  f(x_n)$. Das ist dein \textbf{Codewort}.
\item Beim Empfänger: Wenn höchstens $n - k$ der Werte verloren oder
  verfälscht sind, kann er das ursprüngliche $f$ rekonstruieren — und
  damit auch die Datensymbole.
\end{enumerate}

\subsection*{Warum funktioniert das?}

Weil das Theorem aus dem vorigen Abschnitt sagt: $k$ korrekte
Auswertungen reichen zur Rekonstruktion. Wenn du also $n$ sendest und
maximal $n - k$ verloren gehen, sind immer noch $k$ übrig. Aus denen
rekonstruierst du $f$ und liest die Koeffizienten als Datensymbole
zurück.

\subsection*{Wie viele Fehler korrigieren?}

Hier wird's etwas subtiler:

\begin{itemize}
\item \textbf{Auslöschungen} (\emph{erasures}, du weißt \emph{wo}
  Werte fehlen): bis zu $n - k$ Auslöschungen sind machbar — das ist
  genau das, was wir gerade gesagt haben.
\item \textbf{Fehler} (\emph{errors}, du weißt \emph{nicht}, welche
  Werte verfälscht sind): nur bis zu $\lfloor (n - k) / 2 \rfloor$.
\end{itemize}

Das ist dieselbe Aufteilung wie bei der Hamming-Distanz von Tag 2:
für die Korrektur \emph{ohne Positionswissen} brauchst du etwa doppelt
so viel Redundanz wie für die Erkennung mit Positionswissen.

\begin{bleistiftuebung}{3}{Ein RS$(5, 3)$-Codewort von Hand}
Wir bauen einen Reed-Solomon-Code mit $k = 3$ Datensymbolen und
$n = 5$ Auswertungen über GF($2^3$). Datenpolynom:
\[
  f(x) = 3 + 5x + 2x^2.
\]
Stützstellen: $1, 2, 3, 4, 5$.

\begin{enumerate}[label=(\alph*), leftmargin=*]
\item Berechne $f(1), f(2), f(3), f(4), f(5)$. Das ist das Codewort.
  (Tipp: $f(2)$ haben wir oben schon ausgerechnet — als
  Plausibilitäts-Check.)

\item Aus 5 Werten dürfen $5 - 3 = 2$ \emph{Auslöschungen} korrigiert
  werden. Behaupte: aus den drei Werten $f(1), f(3), f(5)$ allein
  könntest du das Datenpolynom rekonstruieren — die Werte $f(2)$ und
  $f(4)$ sind \emph{verloren}, aber du weißt das.

  Beschreibe in eigenen Worten (kein Rechnen nötig, nur das Prinzip),
  wie der Empfänger vorgeht.

\item Wie viele \emph{nicht-lokalisierte Fehler} (also Verfälschungen
  ohne Positionswissen) kann RS$(5, 3)$ korrigieren?
\end{enumerate}
\end{bleistiftuebung}

\subsection*{Brücke zu früher Gelerntem}

Diese RS-Idee ist die direkte Verallgemeinerung von zwei Sachen, die
wir schon kennen:

\begin{itemize}
\item \textbf{ISBN-10-Detektivarbeit (Tag 2):} dort hast du eine
  einzelne fehlende Ziffer aus den anderen rekonstruiert — sofern du
  wusstest, \emph{welche} fehlt. Genau das ist Auslöschungs-Decodierung
  in $\mathbb{Z}/11\mathbb{Z}$.
\item \textbf{Hamming-Code (Tag 3):} dort haben wir zusätzliche
  Prüfbits an cleveren Positionen eingestreut, sodass die Prüfsummen
  die \emph{Position} eines Fehlers ergaben. Bei RS machen wir es
  anders — wir sampeln das Daten-Polynom an mehr Stellen, als nötig
  wäre, und holen die Position über die mathematische Konsistenz
  zurück.
\end{itemize}

% =============================================================
\section{Python-Werkstatt: Polynom-Klasse und ein RS-Encoder \normalfont(\textit{$\approx$ 35 min})}
% =============================================================

Wir bauen jetzt eine \texttt{Polynom}-Klasse über GF($2^3$). Sie
nutzt die \texttt{GF8}-Klasse aus Tag 5 (lege beide Dateien in
denselben Ordner).

\begin{pythoneinheit}{1}{Polynom-Klasse über GF($2^3$)}
Lege eine neue Datei \texttt{polynom\_gf8.py} an:

\pythoncodeteil{tag6/polynom_gf8.py}{klasse}

Die Methode \texttt{auswerten} fehlt noch — wir benutzen das
\textbf{Horner-Schema}: das ist die effiziente Form, ein Polynom
auszurechnen, ohne mehrfach Potenzen zu bilden. Schau es dir genau
an:

\pythoncodeteil{tag6/polynom_gf8.py}{auswerten}

\paragraph{Aufgabe 1.1 — Verstehen.}
Erkläre kurz (zwei, drei Sätze), warum das Horner-Schema
funktioniert. Tipp: schreibe $a_0 + a_1 x + a_2 x^2 = a_0 + x(a_1 +
x \cdot a_2)$ aus.

\paragraph{Aufgabe 1.2 — Mit der Bleistiftrechnung vergleichen.}
Hast du Bleistiftübung 3 (a) ausgerechnet? Vergleiche dein Ergebnis
mit dem, was Python ausspuckt. Wenn etwas abweicht: liegt der Fehler
auf Papier oder im Code? Beim Debuggen lohnt es sich, ein konkretes
\texttt{f(2)} parallel von Hand und in Python zu rechnen.
\end{pythoneinheit}

\begin{pythoneinheit}{2}{RS-Encoder durch Auswertung}
Mit der \texttt{Polynom}-Klasse und der \texttt{auswerten}-Methode
ist ein Reed-Solomon-Encoder fast schon trivial:

\pythoncodeteil{tag6/polynom_gf8.py}{rs-encode}

\paragraph{Aufgabe 2.1 — Andere Daten.}
Verändere \texttt{daten} und \texttt{stuetzstellen} und schau, wie
sich das Codewort ändert. Was passiert, wenn du eine Stützstelle
zweimal angibst? (Mathematisch: das ist eine schlechte Idee — warum?)

\paragraph{Aufgabe 2.2 — Fehler einbauen.}
Nimm das Codewort, kippe ein Symbol (etwa: ersetze \texttt{codewort[2]}
durch ein anderes GF8-Element) und überlege: welche Information
fehlt dem Empfänger jetzt, um das Datenpolynom zurückzubekommen?
Welcher der beiden Decoder-Fälle aus dem Bleistiftblock ist das —
Auslöschung oder echter Fehler?
\end{pythoneinheit}

% =============================================================
\section{Reflexion und Brücke zu Tag 7 \normalfont(\textit{$\approx$ 15 min})}
% =============================================================

\subsection*{Was wir heute gelernt haben}

\begin{itemize}
\item Polynome funktionieren über jedem Körper — auch über GF($2^3$).
  Addition, Multiplikation, Auswertung sind genau wie in der Schule,
  nur mit anderer Arithmetik in den Koeffizienten.
\item \textbf{Stützstellen-Eindeutigkeit:} ein Polynom Grad $< k$ ist
  durch $k$ Auswertungen eindeutig bestimmt.
\item \textbf{Reed-Solomon} (in seiner einfachsten Form): Daten als
  Polynom, $n > k$ Auswertungen verschicken, Empfänger braucht nur
  irgendwelche $k$ davon.
\item Auslöschungen und Fehler sind unterschiedlich teuer:
  $n - k$ Auslöschungen vs.\ $\lfloor (n - k) / 2 \rfloor$ Fehler.
\end{itemize}

\subsection*{Drei Fragen zum Mitnehmen}

\begin{enumerate}
\item Wir haben heute mit GF($2^3$) und maximal 8 Stützstellen
  gearbeitet. Ein Datamatrix-Code rechnet in GF($2^8$) mit 256
  Elementen. Welche Konsequenz hat die Körpergröße für die maximale
  Codewort-Länge?

\item RS-Encoder durch Auswertung ist konzeptionell glasklar. In der
  Praxis nutzt man oft eine \emph{andere} Form: das Datenpolynom wird
  mit einem festen \textbf{Generator-Polynom} multipliziert, sodass
  das Codewort durch das Generator-Polynom teilbar ist (analog zu CRC
  von Tag 4). Ahnst du, warum der Encoder dann einfacher wird?

\item Beim Decoder hast du zwei Aufgaben: \emph{wo} sind die Fehler,
  und \emph{wie groß} sind sie. Welche der beiden ist die
  schwierigere? (Spoiler für Tag 8.)
\end{enumerate}

\subsection*{Vorschau Tag 7}

Morgen bauen wir den \textbf{Reed-Solomon-Encoder in seiner
gebräuchlichen Form}: mit Generator-Polynom und systematischer
Codierung. „Systematisch“ heißt: das Codewort enthält die
Datensymbole \emph{wörtlich}, nur ein paar Prüfsymbole werden
angehängt — sehr ähnlich zum CRC-Verfahren von Tag 4. Damit ist die
Encoder-Seite für ein Datamatrix-Symbol abgehakt. Den Decoder schauen
wir uns dann an Tag 8 an.
